% Paper 9 (TN Holdings), DRAFT 1, home, 2026-11-14 (owner 21:19 JST "進めて"; scope 21:2x JST: the yardstick in paper
% 9, the light-cone work in paper 10). Built from notes/paper9/part_a_from_paper8.tex, part_a_depth_draft.tex,
% part_b_draft.tex (workplace, numbers checked there against each record's scores.json), the r-thread
% (notes/paper9/r_thread_2026-11-13), the twin test (results/twin-2026-11-13), the yardstick
% (notes/paper9_yardstick) and the XX size test (notes/xxsize_2026-11-14). Every number is quoted from the committed
% record named beside it. Figures: still those of paper 8 (fig3) and the r-thread; to be regenerated for this paper.
\documentclass[10pt,a4paper]{article}
\usepackage[margin=2.2cm]{geometry}
\usepackage{amsmath,amssymb}
\usepackage{booktabs}
\usepackage{graphicx}
\usepackage[hidelinks]{hyperref}
\usepackage{url}
\setlength{\parskip}{4pt}
\setlength{\parindent}{0pt}
\sloppy
\newcommand{\pending}[1]{\textbf{[PENDING: #1]}}

\title{Reading Raw Quantum Dynamics by the Share Kept:\\ IBM's 2023 Utility Observables, the 1D Fermi--Hubbard Model
to 120 Qubits, and One Yardstick for Every Result}
\author{TN Holdings\\ \small contact: \texttt{love.os.architect@proton.me}}
\date{\pending{date of release} (draft 1: 14 November 2026)}

\begin{document}
\maketitle

\begin{abstract}
Papers 7 and 8 of this series ran user-written circuits on IBM Heron devices through an operating layer (psf-os) that
places a circuit from the coupling graph and the failed elements only, adds its own controls to each job and applies
no error mitigation; a result \emph{reads} when its mean squared error is at most half the mean square of the exact
answer. This paper carries the policy to two targets named in the literature and asks what sets where reading stops.
Four of the observables of IBM's 2023 utility experiment, run raw on the published circuit embedded in ibm\_fez,
read up to 8 steps for $\langle Z_{106}\rangle$ and 6 for the weight-10 product, and not for the weight-17 product. The
one-dimensional Fermi--Hubbard model at $U = 4$, with a free-fermion ruler of equal gate count in the same job, read
at every size from 40 to 120 qubits at 2 and 4 Trotter steps and at neither size run at 6. Across both, and across
every other result of the series with an exact answer (179 results in all), we split the error exactly as
$\mathrm{MSE}/\overline{e^2} = (1-r)^2 + \text{rest}$, with $r = \sum ze/\sum e^2$ the share of the exact pattern the
device keeps. Since the rest is never negative, $r < 1 - 1/\sqrt2 \approx 0.293$ means ``does not read'' by
construction; the measured statement is the converse: of the 166 results with $r \ge 0.293$, 160 read. The six that
did not are Hubbard at 6 steps, where the error that is not a shrink (a wrong shape) is as large as the shrink, and
three XX-chain circuits at the line itself. A Clifford twin of the compiled circuit, with the same pulses and an exact
answer at any size, tracked $r$ within 0.051 on a device. A declared test of whether $r$ depends on CZ layers alone
failed: at equal layers, $r$ was lower in an 81-qubit chain than in a 33-qubit one on two devices. No quantum
advantage is claimed: every answer is computed exactly on a classical machine.
\end{abstract}

\section{Introduction}\label{sec:intro}
Paper 7 (doi:10.5281/zenodo.24046673) and paper 8 (doi:10.5281/zenodo.24136512) reported raw device results under
one policy: placement from the coupling graph and the failed elements only, the device chosen by its queue, the job's
own controls (floor, mirror, ring sentinel) in the same submission, no mitigation, dynamical decoupling and twirling
off, answers sealed by hash before each run, every declaration pushed to git before the device record of its run was
fetched. This paper keeps that policy and asks two questions.
\begin{itemize}
\item Does it carry to two targets the field uses as references: IBM's 2023 utility observables (Kim et al.\ 2023), and
an interacting fermion model, the 1D Fermi--Hubbard model, at sizes past exact diagonalisation but inside an exact
classical check?
\item What sets where reading stops? Within a family the boundary lies along depth (Hubbard: every size read at 4
steps, none at 6), but different families stop at very different depths. We look for one quantity that tells, for
every result with an exact answer, whether it reads.
\end{itemize}
Not claimed: quantum advantage; that any mitigation does or does not work (none is used); anything about devices,
days or families not run. Most results are one device and one day; each section says which.

\section{Measures}\label{sec:measures}
For a circuit with exact values $e_i$ and raw device values $z_i$ (expectation values over its measured qubits or
points), the declared verdict is
\[ \textsc{reads} \iff \mathrm{MSE} = \overline{(z-e)^2} \le \tfrac12\,\overline{e^2}, \]
i.e.\ the device's error is at most half of what the all-zero guess would make. Beside it we report
\[ r = \frac{\sum_i z_i e_i}{\sum_i e_i^2}, \qquad
   \frac{\mathrm{MSE}}{\overline{e^2}} = (1 - r)^2 + \underbrace{\frac{\overline{(z - r e)^2}}{\overline{e^2}}}_{\text{rest}}. \]
The split is an identity ($r$ is the least-squares scale, so the cross term vanishes). $r$ is the share of the exact
pattern the device keeps; the rest is everything that is not a uniform shrink (a wrong shape, offsets, shot noise).
Two consequences follow without data. If the device only shrank the signal ($z = re$), \textsc{reads} would hold exactly
when $r \ge 1 - 1/\sqrt2 \approx 0.293$, whatever the size of the signal. Because the rest is never negative,
$r < 0.293$ always means ``does not read''. What a device record can show is how often $r \ge 0.293$ is enough.

\section{IBM's 2023 utility observables, raw}\label{sec:utility}
\subsection{The run of 10 November}
\texttt{results/utility-2026-11-10}. Kim et al.\ (2023) ran a kicked Ising circuit on ibm\_kyiv (Eagle, 127 qubits)
and reported error-mitigated values; here four of their observables were run raw, with the published circuit on the
published graph. Eagle's graph embeds in Heron r2's in two ways. All $R_{ZZ}$ commute, so after $t$ steps an
observable depends only on the gates within $t$ edges of its support; per observable, the first embedding whose light
cone holds no failed element was used, the gates on failed elements (outside the cone) were left out, and the
observable's exact value is the published circuit's. Exact values came from a tensor-network contraction of the
light-cone network, computed twice independently (agreement $2\times10^{-15}$) and checked against the stabilizer
method on all 127 qubits at $\theta_h = 0$ and $\pi/2$; they were sealed before the run.

On ibm\_fez (ibm\_marrakesh had no clear cone for any observable), 17 values of $\theta_h$, 4000 shots:
$\langle Z_{106}\rangle$ after 4 steps (564 CZ) reads; the weight-10 product after 5 steps (705 CZ) reads and keeps 0.61
of its exact $+1$ at $\theta_h = \pi/2$; the weight-17 product and the Fig.~4a product (690 CZ) do not read: at
$\theta_h = \pi/2$ they keep their sign and 7 to 9\% of the exact $-1$. Each product's own mirror falls the same way
(parity 0.50, 0.12, about 0). No comparison is made with IBM's mitigated values.

\subsection{How deep the observables read (11 November)}\label{sec:utility-depth}
\texttt{results/utility-depth-2026-11-11}, ibm\_fez (ibm\_marrakesh placed no clear cone),
prereg/utility\_depth\_2026-11-11.md. The exact values were extended past the 11-10 depths by the same light-cone
contraction (checked at the Clifford points against the stabilizer method on all 127 qubits) and sealed before the
run. One job of 137 circuits, 17 values of $\theta_h$ per observable:

\begin{center}\small
\begin{tabular}{lrrrcc}
\toprule
observable & steps & CZ & light cone (qubits) & MSE / mean square & AGREES (of 17) \\
\midrule
$Z_{106}$ & 4 & 564 & 17 & 0.0198 / 0.4145 reads & 7 \\
$Z_{106}$ & 5 & 705 & 25 & 0.0149 / 0.4073 reads & 7 \\
$Z_{106}$ & 6 & 846 & 31 & 0.0117 / 0.4005 reads & 8 \\
$Z_{106}$ & 7 & 987 & 38 & 0.0110 / 0.3968 reads & 8 \\
$Z_{106}$ & 8 & 1128 & 43 & 0.0108 / 0.3921 reads & 8 \\
weight-10 & 5 & 705 & 37 & 0.0238 / 0.1672 reads & 11 \\
weight-10 & 6 & 846 & 42 & 0.0005 / 0.0038 reads & 17 \\
\bottomrule
\end{tabular}
\end{center}

The prediction that weight-10 would not read at 6 steps failed. Both readings need care. The exact $Z_{106}$ curve
changes little with depth (the 4- and 8-step curves differ by at most 0.11), and its MSE falls with depth because the
device's tail at large $\theta_h$ shrinks: this is reading a slowly changing answer. Weight-10 at 6 steps is a small
signal (mean square 0.0038): the verdict is given as declared. Against 11-10 on the same device and circuits,
$Z_{106}$ at 4 steps had MSE 0.0198 (0.0207) and weight-10 at 5 steps 0.0238 (0.0226).

\section{The one-dimensional Fermi--Hubbard model, raw}\label{sec:hubbard}
\subsection{Circuit, ruler and classical check}
$L$ sites on $2L$ qubits in the interleaved Jordan--Wigner order, N\'eel start, $J = 1$, first-order Trotter steps of
$dt = 0.25$: onsite $R_{ZZ}(U\,dt/2)$ and $R_Z(-U\,dt/2)$, then even and odd bonds by fermionic swaps and $XX+YY$
hopping; checked against the fermion operators ($5.6\times10^{-16}$ at $L = 3$); $10L-8$ CZ and 12 CZ layers per step
on a line, no routing. At $U = 0$ the circuit is free fermions, exact at any size from a $2L\times2L$ single-particle
matrix: the \emph{ruler}, run in the same job on the same chain, with the same CZ count and depth as $U = 4$ (each
onsite $R_{ZZ}$ written as two cancelling CZ). The $U = 4$ answers come from matrix-product states; before any device
run a check with a criterion fixed beforehand covered $L = 20$, 24, 30 at 4, 6, 8 steps (untruncated against truncated
MPS, $2.5\times10^{-8}$ and $7.8\times10^{-8}$; at 8 steps bond dimensions 1024 and 2048 within $8\times10^{-10}$).

\subsection{40 to 120 qubits}
\texttt{results/hubbard-2026-11-11} and \texttt{hubbard-big-2026-11-11}, ibm\_fez, 4000 shots, answers sealed:

\begin{center}\small
\begin{tabular}{rrrrcc}
\toprule
$L$ (qubits) & steps & CZ & CZ per qubit & $U=4$: MSE / mean square & ruler: MSE / mean square \\
\midrule
20 (40) & 2 & 384 & 9.6 & 0.0094 / 0.1814 reads & 0.0063 / 0.0824 reads \\
20 (40) & 4 & 768 & 19.2 & 0.0080 / 0.0460 reads & 0.0395 / 0.1316 reads \\
20 (40) & 6 & 1152 & 28.8 & 0.0141 / 0.0219 does not read & 0.0480 / 0.0512 does not read \\
30 (60) & 2 & 584 & 9.7 & 0.0093 / 0.1725 reads & 0.0043 / 0.0742 reads \\
30 (60) & 4 & 1168 & 19.5 & 0.0068 / 0.0375 reads & 0.0422 / 0.1312 reads \\
30 (60) & 6 & 1752 & 29.2 & 0.0125 / 0.0176 does not read & 0.0361 / 0.0452 does not read \\
45 (90) & 2 & 884 & 9.8 & 0.0191 / 0.1672 reads & 0.0110 / 0.0697 reads \\
45 (90) & 4 & 1768 & 19.6 & 0.0101 / 0.0325 reads & 0.0324 / 0.1383 reads \\
60 (120) & 2 & 1184 & 9.9 & 0.0132 / 0.1637 reads & 0.0062 / 0.0661 reads \\
60 (120) & 4 & 2368 & 19.7 & 0.0073 / 0.0290 reads & 0.0307 / 0.1308 reads \\
\bottomrule
\end{tabular}
\end{center}

Every size from 40 to 120 qubits read at 2 and 4 steps (48 CZ layers at 4); neither size run at 6 steps (72 layers)
read. The predictions that $U = 4$ would read at 6 steps and that the ruler would keep the same share as $U = 4$ (ratio
$r(U4)/r(U0)$ within 0.9--1.1) failed: the ratio held at 2 steps (0.90--0.94) and not at 4 or 6 (1.00--3.40). The
ruler's exact pattern spreads ballistically while the $U = 4$ N\'eel pattern decays slowly, so the two arms carry
different signals; the ruler stayed beside the verdicts and corrected nothing. The exactly-zero site charge
$\langle Z_{i\uparrow}\rangle + \langle Z_{i\downarrow}\rangle$ read 0.05--0.10 on average.

\subsection{One time, two step sizes}
\texttt{results/hub-dt-2026-11-12}, ibm\_fez. At $L = 20$, the same time in two step sizes, each arm against its own
circuit's answers (all four read) and, as declared, against the continuum limit: at $Jt = 1$ the 2-step circuit was
closer than the 4-step one by 0.0009 in MSE (1.7 shot standard errors); at $Jt = 0.5$ the 2-step circuit was closer
than the 1-step one by 0.0172 (17 standard errors), the 1-step circuit's Trotter error (0.0164) being larger than the
2-step circuit's whole device error.

\section{One yardstick for every result}\label{sec:yardstick}
\subsection{The table}
\texttt{notes/paper9\_yardstick/line\_all.py} collects every recorded result of the series with an exact answer and a
declared verdict (written after all were scored): XX chains (75 circuits: the 91-qubit depth scan of 11-12, the size
test and its repeat of 11-14, the twin test of 11-13), Hubbard (31: sections~\ref{sec:hubbard}, the twin test),
IBM 2023 utility (11 observable-depth units, per $\theta_h$ points), 3D kicked Ising (48 arms of 11-10, paper 8), the
daily record (11 problems, 2026-11-05 to 11-14) and the dual-unitary curves at 51, 71 and 91 qubits (3, papers 7 and 8).

\begin{center}\small
\begin{tabular}{lrrr}
\toprule
family & results & read & $r \ge 0.293$ gives the verdict \\
\midrule
XX chain & 75 & 63 & 72 \\
Hubbard (both arms, two step sizes) & 31 & 26 & 28 \\
IBM 2023 utility & 11 & 9 & 11 \\
3D kicked Ising & 48 & 48 & 48 \\
daily record (P4 chain) & 11 & 11 & 11 \\
dual-unitary curves & 3 & 3 & 3 \\
\midrule
all & 179 & 160 & 173 \\
\bottomrule
\end{tabular}
\end{center}

Of the 13 results with $r < 0.293$, none read (the identity). Of the 166 with $r \ge 0.293$, 160 read. The six that do
not:
\begin{itemize}
\item Hubbard at 6 steps, three times ($L = 20$ and 30 on 11-11; $L = 20$ on 11-13): $r$ = 0.38--0.51, but the rest is
0.27--0.32, as large as the shrink. The device keeps much of the pattern and gets its shape wrong.
\item The XX chain at the line itself, three times (MSE / mean square 0.504--0.518, rest 0.03--0.06).
\end{itemize}
Where the error is mostly a shrink (the XX chain, the utility products of weight 10 and 17, the 3D Ising arms), $r$ alone
tells the verdict. Where half the error is rest at every depth (Hubbard at $U = 4$: rest 0.49--0.52 of the ratio, from
the r-thread notes), the line in $r$ sits higher than 0.293 and the boundary comes earlier than $r$ says.

\subsection{Seeing $r$ before the verdict: the Clifford twin}
\texttt{results/twin-2026-11-13}, ibm\_marrakesh, prereg/twin\_2026-11-13.md. On Heron, $R_Z$ is a frame change with no
pulse; rounding every $R_Z$ of a compiled circuit to a multiple of $\pi/2$ gives a circuit with the same pulses (same
CZ, same $\sqrt X$, same timing) whose exact answer a stabilizer simulation gives at any size (psf-os, \texttt{twin}).
Run beside the user's circuit in the same job, its own $r$ is a direct, mitigation-free reading of how much signal the
device keeps on that pulse sequence. On the XX chain at 91 qubits (4, 8, 12 steps) and Hubbard at $L = 20$ (2, 4, 6
steps): $r_{\text{twin}} - r$ = +0.017, +0.018, +0.007 and $-0.051$, $-0.047$, +0.039. The declared bound 0.05 failed
narrowly (Hubbard at 2 steps). ``$r_{\text{twin}} \ge 0.293$'' matched the user circuit's verdict in 5 of 6; the miss is
Hubbard at 6 steps, the same shape error as above, which a twin cannot see.

\section{A limit: $r$ is not set by CZ layers alone}\label{sec:size}
\texttt{notes/xxsize\_2026-11-14}, prereg/xxsize\_2026-11-14.md. On 11-12 (ibm\_fez, 91 qubits) $r$ was nearly the same at
$\theta = \pi/4$ and $\pi/2$ at every depth, which suggested that $r$ depends on the CZ layers a qubit goes through and
not on the circuit's size or signal. A declared test on ibm\_fez and ibm\_marrakesh, XX chains of 33 and 81 qubits,
both angles, 8 to 64 CZ layers, failed on that point and held on the yardstick:
\begin{itemize}
\item angle: $|r(\pi/4) - r(\pi/2)| \le 0.05$ in 12 of 20 pairs (largest 0.074) --- failed;
\item size: $|r(81) - r(33)| \le 0.05$ in 7 of 20 pairs; $r$ lower at 81 qubits in all 20, by up to 0.18 on ibm\_fez ---
failed;
\item ``$r \ge 0.293$'' matched the verdict in 20 of 20 --- held. Readable depth at $\pi/4$: ibm\_marrakesh 64 layers at both
sizes; ibm\_fez 64 at 33 qubits and 32 at 81.
\end{itemize}
On ibm\_marrakesh the 33-qubit chain was exactly the first 33 qubits of the 81-qubit chain, so the same qubits can be
compared inside the two circuits: they keep less inside the longer one, and the difference opens next to where the
longer chain continues and moves inward with depth; a repeat of the same two jobs gave the same order. A stabilizer
simulation of the recorded $\theta = \pi/2$ circuits with uniform local noise gives no such difference ($\pm0.004$), so
the difference needs elements that differ along the paths the dynamics takes. That work, with more devices and days, is
left to paper 10; here it bounds what $r$ is: the share kept on \emph{this} circuit, on \emph{these} qubits, today ---
not a property of the depth alone.

\section{Limits}\label{sec:limits}
\begin{itemize}
\item Most results are one device and one day (utility and Hubbard: ibm\_fez; twin: ibm\_marrakesh).
\item The yardstick table was assembled after the results; only the 20 cells of section~\ref{sec:size} tested the line
under a declaration. Families whose verdicts are on other observables (Ohm's $M$, Schwinger, the 3D light cone) are not
on it.
\item The utility observables at depth include a slowly changing answer ($Z_{106}$) and a small signal (weight-10 at 6
steps); the verdicts are as declared, with that said.
\item The Trotter error is part of each circuit's answer: the checks are of the circuits, not of the continuum models,
except where section~\ref{sec:hubbard} compares with the continuum.
\item The classical checks are the cost of named tools on named machines, not hardness bounds.
\end{itemize}

\section{Reproduction}\label{sec:repro}
Every device result re-scores from its committed raw counts with \texttt{python reproduce.py} in
\texttt{psf-os-lab-notebook} (with psf-os 0.1.1, doi:10.5281/zenodo.24211405, installed): \pending{count of entries for
this paper; all 34 of the notebook re-scored the same from fresh clones on 2026-11-14, before the size test's entries
were added}. The yardstick: \texttt{python notes/paper9\_yardstick/line\_all.py}; the twin: psf-os \texttt{twin}; the
uniform-noise check: \texttt{notes/xxcone\_2026-11-14/stim\_noise.py}.

\section{Conclusion}
Raw, with placement from the coupling graph alone, IBM's 2023 utility observables read to 8 steps ($Z_{106}$) and
6 (weight-10), and the 1D Fermi--Hubbard model read from 40 to 120 qubits at 4 Trotter steps. Across 179 results of
this series, the share of the exact pattern the device keeps, $r$, told the verdict in 173; below the line
$r = 0.293$ nothing reads by construction, and above it 160 of 166 read, the exceptions being a wrong shape (Hubbard at
6 steps) or a result at the line itself. A Clifford twin with the same pulses reads $r$ in the same job without
mitigation. And $r$ belongs to the circuit on its qubits, not to its depth: a longer chain of the same depth kept less.

\section*{Data and code}
Records, declarations and raw counts: \url{https://github.com/TN-Holdings-LLC/psf-os-lab-notebook}; the operating
layer: \url{https://github.com/TN-Holdings-LLC/psf-os} (doi:10.5281/zenodo.24211404); the compiler:
\url{https://github.com/TN-Holdings-LLC/psf-zero}.

\section*{References}
\begin{itemize}
\item Kim, Y.\ et al.\ Evidence for the utility of quantum computing before fault tolerance. Nature 618, 500--505 (2023).
\item Tindall, J., Fishman, M., Stoudenmire, E.\ M., Sels, D. Efficient tensor network simulation of IBM's Eagle kicked
Ising experiment. PRX Quantum 5, 010308 (2024); arXiv:2306.14887.
\item Arute, F.\ et al.\ (Google Quantum AI). Observation of separated dynamics of charge and spin in the Fermi--Hubbard
model. arXiv:2010.07965 (2020).
\item Phasecraft and Google Quantum AI. 2D Fermi--Hubbard dynamics on a superconducting processor. arXiv:2510.26845
\pending{check title and authors}.
\item Phasecraft and Quantinuum. Fermi--Hubbard dynamics on H2. arXiv:2510.26300 \pending{check title and authors}.
\item TN Holdings. Papers 7 and 8 of this series: doi:10.5281/zenodo.24046673, doi:10.5281/zenodo.24136512.
\end{itemize}

\end{document}
